
\medskip
\phantomsection\label{par:sga3-VIA-4.2}%
\noindent\textbf{4.2.---}
Take \(G^\alpha\) in particular to be the connected component
\(G^0\) of the identity in \(G\). Then \(\mathscr S(G^0)\) plainly
contains the image of \(F\) under \(u\), which is precisely the
equivalence class of the identity. On the other hand, if \(F^\beta\)
is a connected component of \(F\), then
\(F^\beta\times_A G^0\) is connected by \hyperref[lem:sga3-VIA-2.1.2]{2.1.2}, so its image under
\(\lambda\) is contained in the connected component of
\(u(F^\beta)\) in \(G\). In other words,
\(\mathscr S(G^0)\) is the union of the connected components that
meet the image of \(F\).

Let us also observe that the open subscheme of \(G\) whose underlying
space is \(\mathscr S(G^0)\) is a subgroup of \(G\); we continue to
denote it by \(\mathscr S(G^0)\). Indeed, inversion on \(G\)
preserves the image of \(F\) and permutes the connected components of
\(G\) that meet this image. It therefore suffices to prove that
\[
  \nu:G\times_A G\longrightarrow G
\]
maps
\(\mathscr S(G^0)\times_A\mathscr S(G^0)\) into
\(\mathscr S(G^0)\). For this we may suppose that \(A\) is an
algebraically closed field: with the notation of \hyperref[par:sga3-VIA-4.1]{4.1},
\(\mathscr S(G^0)\otimes_A\bar k\) is the saturation of
\((G\otimes_A\bar k)^0\) for the equivalence relation defined by the
homomorphism \(u\otimes_A\bar k\).

If \(G^\gamma\) and \(G^\delta\) are connected components of
\(\mathscr S(G^0)\), then \(G^\gamma\times_A G^\delta\) is connected
and its image under \(\nu\) meets the image of \(F\). Consequently,
\[
  \nu(G^\gamma\times_A G^\delta)
\]
is contained in a connected component of \(G\) that meets
\(u(F)\).\begingroup
\renewcommand{\thefootnote}{\fnsymbol{footnote}}%
\footnote[1]{Translator's note: the French source prints
\(u(G^\gamma\times_A G^\delta)\) in this final sentence. The map
defined and used throughout the paragraph is the multiplication
\(\nu:G\times_A G\to G\); the target therefore uses \(\nu\).}%
\endgroup
